\scanpage{5}
\inserted{\unsure{dans 1°}} au cas d'un groupe commutatif, \deleted{\unsure{p}} et par dévissage encore au cas où ce groupe \deleted{\unsure{vérifie}} soit $pK=K$, soit $pK=0$, i.e. les cas 3) $K$ une \unsure{VA.|V.A.}

4) $K$ est unipotent comm. connexe annulé par $p$.

Pour le cas 2), i.e. $K$ fini étale sur $k$, pas néc. commutatif, on note que l'on a une suite exacte (où $\bar{k}$ est la clôture \deleted{Algébrique} séparable de $k$)
\diagram{\begin{tikzcd}[ampersand replacement=\&]
e \arrow[r] \& K(\bar{k}) \arrow[r] \& \unsure{G|C}(\bar{k}) \arrow[r,dash] \& G'(\bar{k}) \arrow[r,dash] \& e
\end{tikzcd}}
et en vertu de l'argument de \unsure{Schreier|Schur}--Serre, on aura
\[
H^1(\pi,G(\bar{k}))-H^1(\pi,G'(\bar{k}))\quad\text{surjectif si}
\]
$H^2(\pi,A)=0$ pour tt $\pi$-module fini $A$. Or cette hypothèse signifie $H^2(k,K)$ pour tt $K$ fini étale sur $k$, et pour un groupe $G$ simple sur $k$, on a
\[
\uline{H^i(k,G)\simeq H^i(\pi,G(\bar{k}))},
\]
\inserted{pour tt $i$, en particulier \unsure{lorsque} $i=\unsure{1|0}$} ce qui prouve le lemme.

\uline{Corollaire} Supposons que $H^2(k,K)=0$ pour tt $K$ fini sur $k$ (pas néc. étale /$k$). Alors les \deleted{conditions} hyp du lemme 1 sont vérifiées, donc aussi la conclusion

Notons d'abord que si $G$ simple sur $k$, alors les $H^i(k,G)$ sont torsion pour $i\unsure{\geq|>}0$ [car $=H^i(\pi,G(\bar{k}))$] donc \deleted{\unsure{si $G$ les|si $G$ est}} \inserted{\unsure{tt}} \unsure{\{|f}\unsure{c}
% archive sheet number: 4

